\SourceLeaf{47}ubi manifesto $\chi$ indicat inclinationem mutuam elementorum $\mathrm{d}t,\mathrm{d}T$, mensuratam per inclinationem normalium utrinque extrorsum respectu spatiorum $s,S$ ductarum, integrationesque per superficies completas utriusque spatii extendi debent.

\Article{9}
Sicuti methodus praecedens divisioni spatii $S$ in elementa prismatica innixa est, ita methodus secunda a divisione eiusdem spatii in elementa pyramidalia petetur. Concipiatur superficies sphaerica radio $=1$ circa centrum $\mu$ descripta atque in elementa infinite parva divisa. Versus punctum talis elementi $\mathrm{d}\Pi$ ducatur a puncto $\mu$ recta, quae superficiem spatii $S$ secet deinceps in punctis $P',P'',P'''$ etc.; distantiae horum punctorum a $\mu$ denotentur per $R',R'',R'''$ etc. Rectae a $\mu$ versus omnia puncta peripheriae elementi $\mathrm{d}\Pi$ ductae formabunt spatium pyramidale, et apud puncta $P',P'',P'''$ etc. e superficie spatii $S$ elementa exsecabunt, quae per $\mathrm{d}T',\mathrm{d}T'',\mathrm{d}T'''$ etc. designamus. Denique sit $Q'$ angulus inter rectam $P'\mu$ atque normalem in elementum $\mathrm{d}T'$ extrorsum ductam, et perinde sint $Q'',Q'''$ etc. inclinationes similium normalium apud puncta $P'',P'''$ etc. ad rectam versus $\mu$ ductam. Ita erit
\[
\mathrm{d}\Pi=\pm\frac{\mathrm{d}T'.\cos Q'}{R'R'}=\mp\frac{\mathrm{d}T''.\cos Q''}{R''R''}=\pm\frac{\mathrm{d}T'''.\cos Q'''}{R'''R'''}\ \text{etc.}
\]
valentibus signis superioribus vel inferioribus, prout $\mu$ est extra vel intra spatium $S$: casus, ubi $\mu$ est in ipsa superficie spatii $S$, adnumerandus est casui priori vel posteriori, prout linea $\mu P'$ extra vel intra spatium $S$ cadit.

Porro patet, pro omnibus partibus spatii $S$ intra illud spatium pyramidale sitis angulum $q$ constantem esse, similique proin modo ut in art. 7 deducimus, si statuatur indefinite
\[
\int\psi r.\mathrm{d}r=-\theta r
\]
constante integrationis ad lubitum accepta, integrale
\[
\int\frac{\mathrm{d}S.\cos q.\psi(\mu,\mathrm{d}S)}{(\mu,\mathrm{d}S)^2}
\]
extensum per omnes partes spatii $S$ intra illud spatium pyramidale sitas, fore in casu priori
\SourceLeaf{48}
\[
=\cos q.\left(\frac{\mathrm{d}T'.\cos Q'.\theta R'}{R'R'}+\frac{\mathrm{d}T''.\cos Q''.\theta R''}{R''R''}+\frac{\mathrm{d}T'''.\cos Q'''.\theta R'''}{R'''R'''}+\ \text{etc.}\right)
\]
in posteriori vero eidem formulae adiiciendum esse terminum $\mathrm{d}\Pi.\cos q.\theta0$.

Iam si haec summatio per omnia superficiei sphaericae elementa colligitur, integrale completum
\[
\int\frac{\mathrm{d}S.\cos q.\psi(\mu,\mathrm{d}S)}{(\mu,\mathrm{d}S)^2}
\]
fiet

I. in casu eo, ubi punctum $\mu$ est extra spatium $S$,
\[
=\int\frac{\mathrm{d}T.\cos q.\cos Q.\theta R}{RR}
\]
denotante $\mathrm{d}T$ indefinite omnia elementa superficiei spatii $S$, atque $Q,R$ illorum respectu eadem, quae antea per literas accentuatas respectu elementorum determinatorum expressa sunt, denique $q$ inclinationem rectae a puncto $\mu$ versus elementum $\mathrm{d}T$ ductae ad rectam nostram fixam.

II. In casu eo, ubi punctum $\mu$ est intra spatium $S$, adiici debet terminus
\[
\theta0.\int\mathrm{d}\Pi.\cos q
\]
ubi $q$ est inclinatio rectae a $\mu$ versus $\mathrm{d}\Pi$ ductae ad rectam fixam, integratioque per totam superficiem sphaericam extendi debet. Sed facile perspicitur, integrale istud, extensum per hemisphaerium id, pro quo $q$ acutus est, fieri $=+\pi$, per hemisphaerium alterum autem $=-\pi$; quapropter integrale completum evanescit, valetque pro hoc casu secundo pure eadem formula, quam pro primo tradidimus. Sed aliter se habet res in casu tertio

III. quoties punctum $\mu$ est in superficie ipsa spatii $S$. Scilicet hic quoque adiiciendus est terminus
\[
\theta0.\int\mathrm{d}\Pi.\cos q
\]
sed integratione per eas tantummodo superficiei sphaericae partes extensa, pro quibus pars initialis rectae a $\mu$ versus $\mathrm{d}\Pi$ ductae cadit intra spatium $S$, sive (siquidem superficies spatii $S$ in puncto $\mu$ neque cuspidem neque aciem offert) pro quibus haec recta facit angulum obtusum cum recta superficiei spatii $S$ in puncto $\mu$ normali extrorsumque ducta. Superest itaque, ut integrale hoc sensu acceptum eruamus.

\SourceLeaf{49}Secent haec normalis atque recta fixa superficiem sphaericam resp. in punctis $G,H$, statuatur arcus $GH=k$, arcus autem inter $G$ et punctum indefinitum superficiei sphaericae $=v$; denique sit $w$ angulus sphaericus inter arcus $k,v$. Ita erit $\cos q=\cos k.\cos v+\sin k.\sin v.\cos w$, et pro $\mathrm{d}\Pi$ accipiendum erit elementum $\sin v.\mathrm{d}v.\mathrm{d}w$. Integrale autem $\int\mathrm{d}\Pi.\cos q$
\[
=\int\int(\cos k.\cos v+\sin k.\sin v.\cos w)\sin v.\mathrm{d}v.\mathrm{d}w
\]
extendi debet a $w=0$ usque ad $w=360^0$, atque a $v=90^0$ usque ad $v=180^0$. Hoc pacto integratio prior suppeditat
\[
\int2\pi\cos k.\cos v.\sin v.\mathrm{d}v
\]
ac dein posterior $-\pi\cos k$.

Ad propositum nostrum hic casus tertius eatenus tantum in considerationem venit, quatenus superficies spatiorum $s,S$ partem quandam finitam communem habent, in qua si punctum $\mu$ reperitur, erit vel $k=0$ vel $=180^0$, adeoque integrale $\int\mathrm{d}\Pi.\cos q$ vel $=-\pi$ vel $=+\pi$, prout scilicet apud punctum $\mu$ spatia $s,S$ sunt vel in eadem plaga vel in plagis oppositis respectu plani utramque superficiem tangentis.

Applicando haec ad integrale nostrum primarium $\int\int\mathrm{d}s.\mathrm{d}S.\varphi(\mathrm{d}s,\mathrm{d}S)$, huius valor fit

I. quoties superficies spatiorum $s,S$ nullam partem finitam communem habent,
\[
=4\pi\sigma\psi0+\int\int\frac{\mathrm{d}t.\mathrm{d}T.\cos q.\cos Q.\theta(\mathrm{d}t,\mathrm{d}T)}{(\mathrm{d}t,\mathrm{d}T)^2}
\]

II. quoties superficies spatiorum $s,S$ partem finitam $=\tau$ communem habent,
\[
=4\pi\sigma\psi0\mp\pi\tau\theta0+\int\int\frac{\mathrm{d}t.\mathrm{d}T.\cos q.\cos Q.\theta(\mathrm{d}t,\mathrm{d}T)}{(\mathrm{d}t,\mathrm{d}T)^2}
\]
ubi signum superius vel inferius valet, prout spatia $s,S$ sunt ab eadem plaga vel a plagis oppositis respectu superficiei communis $\tau$.

III. Quoties superficies spatiorum $s,S$ plures partes finitas discretas communes habent, sit $\tau$ summa earum, quibus spatia $s,S$ ab eadem plaga adiacent, $\tau'$ summa earum, quibus haec spatia a plagis oppositis contigua sunt, eritque integrale nostrum
\SourceLeaf{50}
\[
=4\pi\sigma\psi0+\pi(\tau'-\tau)\theta0+\int\int\frac{\mathrm{d}t.\mathrm{d}T.\cos q.\cos Q.\theta(\mathrm{d}t,\mathrm{d}T)}{(\mathrm{d}t,\mathrm{d}T)^2}
\]
Haec tertia formula omnes casus complecti censeri potest. Integrale duplex per omnia elementa utriusque superficiei extendi debet, denotantque $q,Q$ angulos, quos facit recta bina elementa $\mathrm{d}t,\mathrm{d}T$ iungens cum normalibus in haec elementa extrorsum ductis, directione illius rectae illinc a $\mathrm{d}t$ versus $\mathrm{d}T$, hinc a $\mathrm{d}T$ versus $\mathrm{d}t$ accepta.

\Article{10}
Duae transformationes integralis $\int\mathrm{d}s.\mathrm{d}S.\varphi(\mathrm{d}s,\mathrm{d}S)$ in artt. 8 et 9 evolutae aequali fere concinnitate se commendant, proposito autem nostro posterior magis accommodata est. Problema generale ulterius reduci nequit, nisi ad suppositiones determinatas vel circa spatia $s,S$, vel circa functionem $\varphi$ descendamus. Et quum functio $\varphi$ originem trahat a functione $f$, disquisitionem ulteriorem iam superstruemus eidem hypothesi, a qua ill. \textsc{Laplace} profectus est, puta vires attractivas moleculares in distantiis insensibilibus tantum sensibiles esse. Cui phrasi quum aliquid vagi inhaereat, quamdiu non assignatur unitas, ante omnia observamus, vim attractivam $fr$, per functionem distantiae $r$ expressam, ut cum gravitate $g$ homogenea evadat, antea per massam aliquam multiplicari debere; iam mens illius suppositionis ea est, ut denotante $M$ massam aliquam, qualis in experimentis occurrere potest, puta quam respectu totius terrae pro nihilo habere licet, $Mfr$ semper maneat insensibilis respectu gravitatis, quamdiu $r$ valorem mensuris nostris sensibilem quantumvis parvum habet, dum nihil impediat, quominus valor ipsius $Mfr$ in distantiis insensibilibus non solum sensibilis fieri, sed adeo, decrescente ipsa $r$, omnes limites superare possit. Haud sane sine admiratione deprehendimus, quam gravia ex hac sola hypothesi, dum ceteroquin lex functionis $fr$ tamquam omnino incognita spectatur, eruere liceat, characterem mathematicum prorsus peculiarem prae se ferentia: dum scilicet rebus sic stantibus praecisionem mathematicam absolutam sibi vindicare nequeunt, tamen tantam certissime praecisionem tuentur, ut per nullum experimentum ulla aberratio a veritate absoluta reperiri possit; quamprimum enim successisset, talem aberrationem ulli mensurationi subiicere, suppositio ipsa cessaret.

\SourceLeaf{51}\Article{11}
Supponere licebit, functionem $fr$ (et perinde functionem $Fr$) attractionem denotare, omissa ea parte, quae ipsi $rr$ reciproce proportionalis phaenomenis astronomicis explicandis inservit; haec enim pars, quaecunque sit figura fluidi et vasis, in quovis puncto insensibilem tantummodo modificationem gravitati afferre valet. Crescente itaque $r$ a valore sensibili in infinitum, $fr$ non modo per se insensibilis erit, sed etiam citius decrescet quam $\frac{1}{rr}$. Hinc facile colligitur, etiam integrale $\int fr.\mathrm{d}r$ a valore quocunque sensibili in infinitum extensum insensibile esse, quapropter constantem integrationis $\int fr.\mathrm{d}r=-\varphi r$ ita acceptam supponemus, ut habeatur $\varphi\infty=0$, sive ut sit $\varphi r$ ipse valor integralis $\int fx.\mathrm{d}x$ ab $x=r$ usque ad $x=\infty$ extensus. Hoc pacto $\varphi r$ pro qualibet distantia $r$ denotat quantitatem positivam, sed insensibilem, quamdiu $r$ sensibilis est; contra pro valore insensibili ipsius $r$ non solum sensibilis esse, sed adeo, continuo decrescente distantia $r$, omnes limites superare poterit, sive secundum vulgarem loquendi modum nihil obstat, quominus sit $\varphi0=\infty$.

\Article{12}
Inde quod functio $\varphi r$ pro quovis valore sensibili ipsius $r$ insensibilis est et crescente $r$ continuo decrescit, statim quidem sequitur, integrale $\int rr\varphi r.\mathrm{d}r$ a valore aliquo sensibili usque ad alium maiorem extensum etiamnum insensibile manere, dummodo posterior sit intra ambitum eorum, circa quos experimenta instituere licet: sed neutiquam ex illa proprietate sola concludere fas esset, integrale insensibile manere, ad quantumvis magnum intervallum integratio extendatur. Calculi ill. \textsc{Laplace} ita quidem pronunciati sunt, ut talem suppositionem involvant; at dum natura functionis $\varphi r$ incognita est, consultius videtur, ab omnibus suppositionibus hypotheticis, quibus supersedere possumus, abstinere. Quum itaque constans integrationis $\int rr\varphi r.\mathrm{d}r=-\psi r$ arbitrio relicta sit, sufficiat nobis, eam ita electam supponere, ut fiat $\psi r=0$ pro valore aliquo sensibili ipsius $r$ arbitrario, sed intra ambitum dimensionum corporum, circa quae experimenta instituere licet. Hoc pacto $\psi r$ pro quovis alio eiusmodi valore semper insensibilis erit (positiva pro minori, negativa pro maiori), sed nihil hinc obstat, quominus pro valore insensibili ipsius $r$ sensibilis evadere possit: addere tamen oportet, phaenomenorum explicationem postulare ut decrescente distantia $r$ in infinitum, valor ipsius $\psi r$ semper maneat finitus, sive ut $\psi0$ sit quantitas finita. Ce%
